\documentclass{article}
\usepackage{booktabs} % For prettier tables
\usepackage{luastats}
\usepackage{tabularx} % For adjustable column widths
\usepackage{amsmath, amssymb} % For mathematical symbols
\usepackage{geometry,url} % For better page margins
\geometry{a4paper, margin=1in} % Adjust margins
\usepackage{multirow}

\title{The \texttt{luastats} Package in LaTeX}
\author{Chetan Shirore , Ajit Kumar\\ and\\ Anupam Maurya}
\date{March 28, 2024}

\begin{document}
	
	\maketitle
	
	\section{Introduction}
	
	The luastats package is a LuaLaTeX-based statistical computation library designed to provide essential statistical functions using Lua scripting within LaTeX. It aims to bridge the gap between LaTeX's typesetting capabilities and Lua's computational power, enabling users to perform statistical calculations directly within their LaTeX documents without relying on external tools or complex macros.
	
	\section{Motivation Behind \underline{luastats}}
	
	Statistical analysis is a fundamental aspect of many academic and research fields, including mathematics, physics, economics, and social sciences. While LaTeX excels at formatting mathematical expressions and reports, its built-in calculation capabilities are limited. Traditionally, users have relied on tools like R, Python, or MATLAB for data analysis before manually integrating results into their LaTeX documents. The luastats package eliminates this extra step by allowing users to compute statistical measures dynamically within their \LaTeX \;workflow using Lua.
	
	\subsection{Key Features}
	\begin{enumerate}
		\item Seamless Lua Integration: Leverages LuaLaTeX's embedded Lua scripting environment to perform calculations efficiently.
		\item Essential Statistical Functions: Provides core statistical operations,including:
		\begin{itemize}
			\item Mean (average) computation
			\item Median calculation
			\item Rounding functions
		\end{itemize}
		\item Dynamic Computation in LaTeX: Users can pass data directly in LaTeX and receive computed results in real time, enhancing automation in document generation.
		\item User-Friendly Syntax: The package is designed to be intuitive, requiring minimal coding knowledge to use its functions.
		\item Efficient Execution: Lua's lightweight nature ensures fast computations, making it a practical alternative to heavier computational tools.
	\end{enumerate}
	
	\section{Who Should Use?}
	\begin{itemize}
		\item Researchers \& Academics: Ideal for those who work with data-driven reports and need real-time statistical processing in LaTeX.
		\item Students \& Educators: Helps in teaching statistics, probability, and mathematical analysis without switching between different software.
		\item Technical Writers: Useful for professionals preparing reports, articles, or books that require embedded calculations.
	\end{itemize}
	
	\section{Installation and License}
	
	To install the \texttt{luastats} package, place the \texttt{luastats.sty} file in your LaTeX directory. Include the package in your document with:
	
	\begin{verbatim}
		\usepackage{luastats}
	\end{verbatim}

	\section{Basic Statistical Method}

\begin{table}[h]
	\centering
	\renewcommand{\arraystretch}{2} % Adjust row spacing
	\begin{tabular}{l|l|l}
		\toprule
		\textbf{Code} & \textbf{Uses} & \textbf{Output} \\
		\midrule
		\texttt{\textbackslash luaMean\{1,2,3,4,5\}} & Computes the arithmetic mean & \luaMean{1,2,3,4,5} \\
		\texttt{\textbackslash luaMedian\{1,2,3,4,5\}} & Finds the median & \luaMedian{1,2,3,4,5} \\
		\texttt{\textbackslash luaMode\{1,2,2,3,3,3,4,4\}} & Determines mode & \luaMode{1,2,2,3,3,3,4,4} \\
		\texttt{\textbackslash luaVariance\{1,2,3,4,5\}} & Computes variance & \luaVariance{1,2,3,4,5} \\
		\texttt{\textbackslash luaSD\{1,2,3,4,5\}} & Standard deviation & \luaSD{1,2,3,4,5} \\
		\texttt{\textbackslash luaRange\{1,2,3,4,5\}} & Range (max - min) & \luaRange{1,2,3,4,5} \\
		\texttt{\textbackslash luaQuartiles\{1,2,3,4,5,6,7,8,9,10\}} & Computes quartiles & \luaQuartiles{1,2,3,4,5,6,7,8,9,10} \\
		\texttt{\textbackslash luaIQR\{1,2,3,4,5,6,7,8,9,10\}} & Interquartile range & \luaIQR{1,2,3,4,5,6,7,8,9,10} \\
		\texttt{\textbackslash luaSkewness\{1,2,3,4,5,6,7,8,9,10\}} & Measures skewness & \luaSkewness{1,2,3,4,5,6,7,8,9,10} \\
		\texttt{\textbackslash luaKurtosis\{1,2,3,4,5,6,7,8,9,10\}} & Measures kurtosis & \luaKurtosis{1,2,3,4,5,6,7,8,9,10} \\
		\texttt{\textbackslash luaCorrelation\{1,2,3,4\}\{2,4,6,8\}} & Pearson correlation & \luaCorrelation{1,2,3,4}{2,4,6,8} \\
		\texttt{\textbackslash luaCovariance\{1,2,3,4\}\{2,4,6,8\}} & Covariance & \luaCovariance{1,2,3,4}{2,4,6,8} \\
		\texttt{\textbackslash luaRegression\{1,2,3,4\}\{2,4,6,8\}} & Linear regression & \luaRegression{1,2,3,4}{2,4,6,8} \\
		\texttt{\textbackslash luaChiSquare\{5,10,15\}\{6,9,15\}} & Chi-square test & \luaChiSquare{5,10,15}{6,9,15} \\
		\texttt{\textbackslash luaEMA\{1,2,3,4,5\}\{0.5\}} & Exponential moving average & \luaEMA{1,2,3,4,5}{0.5} \\
		\texttt{\textbackslash luaWeightedMean\{1,2,3,4\}\{0.1,0.2,0.3,0.4\}} & Weighted mean & \luaWeightedMean{1,2,3,4}{0.1,0.2,0.3,0.4} \\
		\texttt{\textbackslash luaZScores\{1,2,3,4,5\}} & Converts to Z-scores & \luaZScores{1,2,3} \\
		\bottomrule
	\end{tabular}
	\caption{Statistical Functions Reference}
	\label{tab:stats_functions}
\end{table}

\texttt{\textbackslash luaCSVsummary\{data.csv\}} Summary from CSV :
\[
\luaCSVsummary{data.csv}
\]

\DataSet{data1}{1,2,3,4,5,6,7,8,9,10}

\begin{table}[h]
	\centering
	\renewcommand{\arraystretch}{2} % Adjust row spacing
	\begin{tabular}{l|l|l}
		\toprule
		\textbf{Code} & \textbf{Uses} & \textbf{Output} \\
		\midrule
		\texttt{\textbackslash DataSet\{data1\}\{1,2,3,4,5\}} & Defines dataset &  -- \\
		\texttt{\textbackslash luastat\{data1\}\{mean\}} & Computes mean & \luastat{data1}{mean} \\
		\texttt{\textbackslash luastat\{data1\}\{median\}} & Computes median & \luastat{data1}{median} \\
		\texttt{\textbackslash luastat\{data1\}\{mode\}} & Computes mode & \luastat{data1}{mode} \\
		\texttt{\textbackslash luastat\{data1\}\{variance\}} & Computes variance & \luastat{data1}{variance} \\
		\texttt{\textbackslash luastat\{data1\}\{standard\_deviation\}} & Computes standard deviation & \luastat{data1}{standard_deviation} \\
		\texttt{\textbackslash luastat\{data1\}\{range\}} & Computes range & \luastat{data1}{range} \\
		\texttt{\textbackslash luastat\{data1\}\{quartiles\}} & Computes quartiles & \scriptsize{\luastat{data1}{quartiles}} \\
		\texttt{\textbackslash luastat\{data1\}\{iqr\}} & Computes interquartile range & \luastat{data1}{iqr} \\
		\texttt{\textbackslash luastat\{data1\}\{skewness\}} & Measures skewness & \luastat{data1}{skewness} \\
		\texttt{\textbackslash luastat\{data1\}\{kurtosis\}} & Measures kurtosis & \luastat{data1}{kurtosis} \\
		\bottomrule
	\end{tabular}
	\caption{Statistical Functions for a Single Dataset}
\end{table}


% \DefineStatAnalysis{analysis1}{1,2,3,4}{2,4,6,8}

% \begin{table}[h]
	%     \centering
	%     \renewcommand{\arraystretch}{1.3} % Adjust row spacing
	%     \begin{tabular}{l|l|l}
		%         \toprule
		%         \textbf{Code} & \textbf{Uses} & \textbf{Output} \\
		%         \midrule
		%         \texttt{\textbackslash DefineStatAnalysis\{analysis1\}\{1,2,3,4\}\{2,4,6,8\}} & Defines two datasets & -- \\
		%         \texttt{\textbackslash StatMethod\{analysis1\}\{correlation\}} & Computes correlation & \StatMethod{analysis1}{correlation} \\
		%         \texttt{\textbackslash StatMethod\{analysis1\}\{covariance\}} & Computes covariance & \StatMethod{analysis1}{covariance} \\
		%         \texttt{\textbackslash StatMethod\{analysis1\}\{linear\_regression\}} & Computes linear regression & \StatMethod{analysis1}{linear_regression} \\
		%         \bottomrule
		%     \end{tabular}
	%     \caption{Statistical Functions for Two Datasets}
	% \end{table}

	\newpage
The package provides a comprehensive set of statistical functions for LaTeX documents. It simplifies the inclusion of PMFs, PDFs, and CDFs for various distributions, making it a valuable tool for researchers, statisticians, and students.
\begin{itemize}
	\item Probability Mass Functions (PMFs)
	\item Probability Density Functions (PDFs)
	\item Cumulative Distribution Functions (CDFs)
\end{itemize}

\begin{table}[ht]
	\centering
	\renewcommand{\arraystretch}{1.5}
	\small % Reduce font size for better fit
	\begin{tabularx}{\textwidth}{lXl}
		\toprule
		\textbf{Name and Syntax} & \textbf{Output} \\
		\midrule
		Binomial \\ \verb|\luabinomialPMF| & $\luabinomialPMF$ \\ \\
		Geometric \\ \verb|\luageometricPMF| & $\luageometricPMF$ \\\\ 
		Negative Binomial \\ \verb|\luanegBinomialPMF| & $\luanegBinomialPMF$ \\\\
		Poisson \\ \verb|\luapoissonPMF| & $\luapoissonPMF$ \\\\
		Hypergeometric \\ \verb|\luahypergeometricPMF| & $\luahypergeometricPMF$ \\\\
		Discrete Uniform \\ \verb|\luadiscreteUniformPMF| & $\luadiscreteUniformPMF$ \\\\
		Multinomial \\ \verb|\luamultinomialPMF| & $\luamultinomialPMF$ \\\\
		Zipf \\ \verb|\luazipfPMF| & $\luazipfPMF$ \\\\
		Negative Hypergeometric \\ \verb|\luanegHypergeometricPMF| & $\luanegHypergeometricPMF$ \\\\
		\bottomrule
	\end{tabularx}
	\caption{Probability Mass Functions (PMFs)}
\end{table}
\newpage

\begin{table}[ht]
	\centering
	\renewcommand{\arraystretch}{1.5}
	\small % Reduce font size for better fit
	\begin{tabularx}{\textwidth}{lXl}
		\toprule
		\textbf{Name} & \textbf{Syntax} & \textbf{Output} \\
		\midrule
		Uniform Distribution & \verb|\luauniPDF| & $\luauniPDF$ \\\\
		Normal (Gaussian) Distribution & \verb|\luanormalPDF| & $\luanormalPDF$ \\\\
		Exponential Distribution & \verb|\luaexpPDF| & $\luaexpPDF$ \\\\
		Gamma Distribution & \verb|\luagammaPDF| & $\luagammaPDF$ \\\\
		Beta Distribution & \verb|\luabetaPDF| & $\luabetaPDF$ \\\\
		Cauchy Distribution & \verb|\luacauchyPDF| & $\luacauchyPDF$ \\\\
		Laplace Distribution & \verb|\lualaplacePDF| & $\lualaplacePDF$ \\\\
		Student's t Distribution & \verb|\luastudentTPDF| & $\luastudentTPDF$ \\\\
		Weibull Distribution & \verb|\luaweibullPDF| & $\luaweibullPDF$ \\\\
		Pareto Distribution & \verb|\luaparetoPDF| & $\luaparetoPDF$ \\\\
		Rayleigh Distribution & \verb|\luarayleighPDF| & $\luarayleighPDF$ \\\\
		Log-Normal Distribution & \verb|\lualognormalPDF| & $\lualognormalPDF$ \\\\
		\bottomrule
	\end{tabularx}
	\caption{Probability Density Functions (PDFs)}
\end{table}
\newpage

\begin{table}[ht]
	\centering
	\renewcommand{\arraystretch}{1.5}
	\small % Reduce font size for better fit
	\begin{tabularx}{\textwidth}{lXl}
		\toprule
		\textbf{Name} & \textbf{Syntax} & \textbf{Output} \\
		\midrule
		Uniform Distribution & \verb|\luauniCDF| & $\luauniCDF$ \\\\
		Normal (Gaussian) Distribution & \verb|\luanormalCDF| & $\luanormalCDF$ \\\\
		Exponential Distribution & \verb|\luaexpCDF| & $\luaexpCDF$ \\\\
		Gamma Distribution & \verb|\luagammaCDF| & $\luagammaCDF$ \\\\
		Beta Distribution & \verb|\luabetaCDF| & $\luabetaCDF$ \\\\
		Cauchy Distribution & \verb|\luacauchyCDF| & $\luacauchyCDF$ \\\\
		Laplace Distribution & \verb|\lualaplaceCDF| & $\lualaplaceCDF$ \\\\
		Student's t Distribution & \verb|\luastudentTCDF| & $\luastudentTCDF$ \\\\
		Weibull Distribution & \verb|\luaweibullCDF| & $\luaweibullCDF$ \\\\
		Pareto Distribution & \verb|\luaparetoCDF| & $\luaparetoCDF$ \\\\
		Rayleigh Distribution & \verb|\luarayleighCDF| & $\luarayleighCDF$ \\\\
		Log-Normal Distribution & \verb|\lualognormalCDF| & $\lualognormalCDF$ \\\\
		\bottomrule
	\end{tabularx}
	\caption{Cumulative Distribution Functions (CDFs)}
\end{table}




	\begin{table}[h]
	\centering
	\renewcommand{\arraystretch}{2} % Adjust row spacing
	\setlength{\tabcolsep}{10pt} % Adjust column spacing
	\begin{tabular}{|c|c|c|}
		\hline
		\textbf{Distribution} & \textbf{Parameters} & \textbf{Generated Values} \\
		\hline
		\multirow{2}{*}{Normal} & $\mu = 0, \sigma = 1, n=5$ & \luarandomnormal{0}{1}{5} \\
		& \verb|\luarandomnormal{0}{1}{5}| & \\
		\hline
		\multirow{2}{*}{Uniform} & $a = 0, b = 1, n=5$ & \luarandomuniform{0}{1}{5} \\
		& \verb|\luarandomuniform{0}{1}{5}| & \\
		\hline
		\multirow{2}{*}{Exponential} & $\lambda = 1, n=5$ & \luarandomexponential{1}{5} \\
		& \verb|\luarandomexponential{1}{5}| & \\
		\hline
		\multirow{2}{*}{Binomial} & $n = 10, p = 0.5, samples=5$ & \luarandombinomial{10}{0.5}{5} \\
		& \verb|\luarandombinomial{10}{0.5}{5}| & \\
		\hline
		\multirow{2}{*}{Poisson} & $\lambda = 3, samples=5$ & \luarandompoisson{3}{5} \\
		& \verb|\luarandompoisson{3}{5}| & \\
		\hline
		\multirow{2}{*}{Chi-Square} & $df = 5, samples=5$ & \luarandomchisquare{5}{5} \\
		& \verb|\luarandomchisquare{5}{5}| & \\
		\hline
		\multirow{2}{*}{t-Distribution} & $df = 10, samples=5$ & \luarandomt{10}{5} \\
		& \verb|\luarandomt{10}{5}| & \\
		\hline
		\multirow{2}{*}{F-Distribution} & $df_1 = 5, df_2 = 10, samples=5$ & \luarandomf{5}{10}{5} \\
		& \verb|\luarandomf{5}{10}{5}| & \\
		\hline
	\end{tabular}
	\caption{Random Number Generation from Various Distributions}
\end{table}
	\section{Plotting Probability Distributions}
This package provides plotting functionality for different distributions using \texttt{TikZ} and \texttt{pgfplots}.

\subsection{Normal Distribution}
\subsubsection{Probability Density Function (PDF)}
\begin{verbatim}
	\luanormpdf{mean}{stddev}{range}{step}
\end{verbatim}
Example:
\begin{center}
	\luanormpdf{0}{1}{-3,3}{0.1}
\end{center}

\subsubsection{Cumulative Distribution Function (CDF)}
\begin{verbatim}
	\luanormcdf{mean}{stddev}{range}{step}
\end{verbatim}
Example:
\begin{center}
	\luanormcdf{0}{1}{-3,3}{0.1}
\end{center}

\subsection{Uniform Distribution}
\subsubsection{Probability Density Function (PDF)}
\begin{verbatim}
	\luaunipdf{a}{b}{range}
\end{verbatim}
Example:
\begin{center}
	\luaunipdf{0}{1}{-1,2}
\end{center}

\subsubsection{Cumulative Distribution Function (CDF)}
\begin{verbatim}
	\luaunicdf{a}{b}{range}
\end{verbatim}
Example:
\begin{center}
	\luaunicdf{0}{1}{-1,2}
\end{center}

\subsection{Exponential Distribution}
\subsubsection{Probability Density Function (PDF)}
\begin{verbatim}
	\luaexppdf{lambda}{range}{step}
\end{verbatim}
Example:
\begin{center}
	\luaexppdf{1}{0,5}{0.1}
\end{center}

\subsubsection{Cumulative Distribution Function (CDF)}
\begin{verbatim}
	\luaexpcdf{lambda}{range}{step}
\end{verbatim}
Example:
\begin{center}
	\luaexpcdf{1}{0,5}{0.1}
\end{center}

\section{Histogram and KDE (Kernel Density Estimation)}
\subsection{Histogram and Smoothed KDE}
\begin{verbatim}
	\luadistplot{data}{bins}{range}{step}
\end{verbatim}

Example:
\begin{center}
	\luadistplot{\luarandomnormal{0}{1}{5000}}{25}{-5,5}{0.5}
\end{center}

	
\end{document}
